\section{Polynomial multiplication}
From a computational standpoint, the most expensive operations in HQC are by far the polynomial multiplication.
In this section we will describe the various methods employed to implement polynomial multiplication in HQC.

	% descrivere la moltiplicazione shift & add semplice
	\subsection{Shift and Add}\label{subsec:shift_add}
		In the third round reference implementation of HQC, the polynomial multiplication is implemented using a variant of the schoolbook
		\textit{shift and add} method, recalling that the multiplication in $\mathcal{R}$ of a polynomial by a monomial of degree $i$ 
		can be implemented as a circular shift to the left (\clshift) 	of $i$ positions. Given $a(x)$ in sparse form and $b(x)$ in dense form, the algorithms multiplies $b(x)$ by $x^w$, 
		for all $w$s, where $w$ is a non-zero coefficient of $a(x)$, and then adds the results together:

		\begin{algorithm2e}[H]
			\caption{Sparse polynomial multiplication}
			\KwIn{$a(x) \in \mathcal{R}_\omega$ in sparse form, $b(x) \in \mathcal{R}$ in dense form}
			\KwOut{$c(x) = a(x) \cdot b(x) \in \mathcal{R}$ in dense form}
			\For{$i \gets 0$ \KwTo \texttt{ts}}{
				$\texttt{table}[i] \gets b(x) \clshift i$\;
			}
			\For{$i \gets 0$ \KwTo $\texttt{wt}(a(x))$}{
				$w \gets a(x)[i]$\;
				$\texttt{dec} \gets w \mod \texttt{ts}$\;
				$\texttt{s} \gets \lfloor \frac{w}{\texttt{ts}} \rfloor$\;
				$\texttt{tmp} \gets x^\texttt{s}\cdot x^\texttt{ts}$\;
				$c(x) \gets c(x) + \texttt{table}[\texttt{dec}] \cdot \texttt{tmp}$\;
			}
			\KwRet $c(x) \mod (x^n+1)$\;
		\end{algorithm2e}
		The idea of this variant is to precompute the first \texttt{ts} shifts of $b(x)$ and store them in a table, where \texttt{ts} is set to 16;
		this allows to slightly improve the performance of the algorithm and to shuffle the memory accesses to \texttt{table}.\\
		Since $a$ is very sparse, this approach is extremely fast; however, the high speed achieved is counterbalanced by the
		non-constant time execution of the algorithm, which is a major drawback for the security of the implementation. 

	% descrivere la versione a tempo costante della mul shift & add 
	\subsection{Constant-time shift and add}\label{subsec:const_shift_add}
		Due to the high sparsity of one of the operands, we would like to exploit the fact that the algorithm is very fast to implement a constant-time version of it;
		such an algorithm is implemented in LEDAcrypt~\cite{baldi2020ledacrypt} and is based on \textit{shifts with conditional commitment}.\\

		\begin{algorithm2e}[H]
			\caption{Constant-time monomial multiplication}
			\KwIn{$a$ exponent of a monomial $x^a \in \mathds{Z}_2[x]$ with $0 \leqslant a < n$, \newline
					$b(x) \in \mathcal{R}$ in dense form}
			\KwOut{$c(x) = x^a \cdot b(x) \in \mathcal{R}$ in dense form}
			\KwData{$m = \lceil \log_2(a) \rceil$ number of bits needed to encode $a$ in binary, \newline
					\textit{i.e.} $a = (\beta_{m-1}, \cdots, \beta_0)_{\texttt{bin}}$}

			$c(x) \gets b(x) \clshift (a \mod \texttt{ws})$\;
			\For{$i \gets \log_2(\texttt{ws})$ \KwTo $(m-1)$} {
				$\texttt{tmp} \gets b(x) \clshift i$\;
				$c(x) \gets c(x) + \texttt{tmp}$\;
				$\texttt{cond} \gets (\beta_{i} == 1)$\;
				$\texttt{mask} \gets 0 - (\texttt{WORD})(\texttt{cond})$\;
				$c(x) \gets c(x) + ((c(x) \wedge \texttt{mask}) \vee (\texttt{tmp}(x) \wedge \texttt{mask}))$\;
			}
			\KwRet $c(x)$\;
		\end{algorithm2e}
		With this algorith we compute the $i$-th shift of $b(x)$ regardless of the value of $\beta_i$, but we sum it to the result only if
		$\beta_i = 1$ (the mask at line 6 contains either all zeros or all ones, '$\vee$' and '$\wedge$' correspond to the bitwise OR and AND respectively).
	% descrivere la moltiplicazione con trasformate

\section{Masking}
    The basic principle of masking to secure a cryptographic algorithm is to split every sensitive value $x$ into $d+1$ \textit{shares},
    in such a way that the following relation holds:
    \begin{equation}
        x = x_0 \oplus x_1 \oplus \cdots \oplus x_d
    \end{equation}
    The shares $x_1, \cdots, x_d$ are called \textit{masks} and are randomly sampled, while the secret $x_0$ is processed in 
    such a way that the relation above holds. The goal of masking is to make the sensitive value indistinguishable from
    a random value, so that the observation of a side channel does not reveal any sensitive information.

    Unlike many other cryptosystems, HQC is \textit{linear} in all its operations, making its primitives to be easily masked entirely.
    In the next algorithms we will use the following notation and conventions:
    \begin{itemize}
        \item $d$: the amount of shares
        \item $\ssplit{a} = (a_0,\ldots,a_{d-1})$: additive share splitting of $a$ in $d$ shares over $\mathcal{R}$
        \item $\sample$: uniform sampling from the right hand side of the equation
        \item $(a+\sum_{i=0}^{d-2} \mathtt{rnd}_i)$: masking of $a$ with $d-1$ masks sampled from the same field of $a$
        \item \textsc{Collapse}(a, \ssplit{b}): sum all shares of $b$, in order, to $a$
    \end{itemize}
    In our implementation we will consider only the case $d=2$.

    \subsection{Keygen}
    The masking of the key generation algorithm is straightforward: after sampling the private key, we mask $y$ and multiply 
    by $h$: the multiplication between a polynomial and a share-split is performed by multiplying each share by $h$, so the result is
    itself share-split.\\
    \begin{algorithm2e}[H]
        \caption{Keygen}
        \KwIn{none}
        \KwOut{$\kpub \in \mathcal{R}\times \mathcal{R}; \kpri \in \mathcal{R} \times \mathcal{R}$}
        $h \sample \mathcal{R}$\;
        $(x,y) \sample \mathcal{R}_\omega \times \mathcal{R}_\omega$\;
        $\ssplit{y} \gets (y+\sum_{i=0}^{d-2} \mathtt{rnd}_i)$\;
        $\ssplit{s} \gets h \ssplit{y}$\;
        $\ssplit{s} \gets x + \ssplit{s}$\tcp*{sum $x$ only to the first share}
        \KwRet $\kpri=(x,y)$, $\kpub = (h,\textsc{Collapse}(\ssplit{s}))$\;
        \end{algorithm2e}

    \subsection{Encrypt}
    In the encryption it is important to protect, other than the message $m$ and the public key $s$, also the random polynomials
    $e,r_1$ and $r_2$, as their leakage would allow an attacker to remove the \textit{excess} errors and decode $m$ from the ciphertext.
    
    \begin{algorithm2e}[H]
        \caption{Encrypt($\kpub$, $m$)}
        \KwData{$\mathbf{G}$: a $\mathds{F}_2^{k \times n}$ matrix generating an efficiently correctable code}
        \KwIn{$\kpub$, $m\in \{0,1\}^k$}
        \KwOut{$c = (u, v) \in \mathcal{R} \times \mathcal{R}$}
        $e \sample \mathcal{R}_{\omega_e}$\;
        $(r_1,r_2) \sample \mathcal{R}_{\omega_r} \times \mathcal{R}_{\omega_r}$\;
        $\ssplit{r_1} \gets (r_1+\sum_{i=0}^{d-2} \mathtt{rnd}_i)$\;
        $\ssplit{r_2} \gets (r_2+\sum_{i=0}^{d-2} \mathtt{rnd}_i)$\;
        $\ssplit{u} \gets \ssplit{r_1} + h\ssplit{r_2}$\;
        $\ssplit{m} \gets (m+\sum_{i=0}^{d-2} \mathtt{rnd}_i) $\;
        $\ssplit{v} \gets \ssplit{m}\mathbf{G} + s \ssplit{r_2} + e $\;
        \KwRet $(\textsc{Collapse}(\ssplit{u}),\textsc{Collapse}(\ssplit{v}))$\;
    \end{algorithm2e}
    When calculating \ssplit{v} it is important to perform the operations reported at line 7 from left to right, to avoid leaking
    information about $e$; in alternative, we can also mask $e$ and then sum it to \ssplit{v}.


    \subsection{Decrypt}
    In the decryption we exploit the linearity of the correctable code: after masking $y$ and multiplying by $u$, we 
    generate a random codeword $\ssplit{\bar{m}}$ and encode it into \ssplit{\bar{c}} to protect $v-uy$ during the decoding.
    After the decoding we can collapse $\ssplit{\bar{m}}$ and the decoded message: since the code is linear, the final result
    is the original message, as if we had performed the decoding on $v-uy$ directly.
    \begin{algorithm2e}
        \caption{Decrypt($\kpri$, $c$)}
        \KwData{$\mathbf{G}$: a $\mathds{F}_2^{k \times n}$ matrix generating an efficiently correctable code}
        \KwIn{$\kpri=(x,y) \in \mathcal{R}_\omega \times \mathcal{R}_\omega$, $c = (u,v) \in \mathcal{R} \times \mathcal{R}$}
        \KwOut{$m\in \{0,1\}^k$}
        $\ssplit{y} \gets (y+\sum_{i=0}^{n-2} \mathtt{rnd}_i)$\;
        $\ssplit{uy} \gets u \ssplit{y}$\;
        $\ssplit{\bar{m}} \sample (\{0,1\}^k)^d$\;
        $\ssplit{\bar{c}} \gets \ssplit{\bar{m}}\mathbf{G}$\; 
        $a \gets \text{\sc Collapse}(v ,\ssplit{\bar{c}})$\;
        $m \gets \text{\sc Decode}(\text{\sc Collapse}(a,\ssplit{uy}))$\;
        \KwRet $\text{\sc Collapse}(m ,\ssplit{\bar{m}})$\;
    \end{algorithm2e}